\section{Introduction}
\label{sec:intro}
Private information retrieval (PIR) hides a client's requested indices
from the server~\cite{chor1998private,kushilevitz1997replication}.
When one task needs several records, repeating a single-record protocol
also repeats query generation, evaluation and extraction. Concurrent
execution can overlap that work, so reducing the number of queries alone
does not establish a completed-task latency benefit.

This paper studies a restricted alternative within the linear evaluation
interface of XPIR~\cite{melchor2016xpir}. The client encrypts one weighted
selector vector for an ordered tuple of distinct indices. Each record
is divided into bounded integer segments and encoded into polynomial
coefficients. Public radix weights place the selected segments in
different digit positions of each reply coefficient, allowing ordinary
digit extraction after decryption. A query is one vector of \(N\)
ciphertexts, not a compressed single ciphertext; the reply contains
one ciphertext per polynomial block.

Encrypted weighted multi-value retrieval is established prior
work~\cite{mulityquery}, and radix representation is standard. The question
here is how to specify its encoding, capacity and integer-noise conditions
for an XPIR-style RLWE interface, and what performance this instantiation
exhibits against concurrent same-profile repeated retrieval. Coefficient
capacity, backend API support and finite functional execution are distinct
constraints. In particular, a formal law cannot be inferred from a
native sampler merely because observed replies decode correctly.

The paper makes three contributions:
\begin{itemize}
\item An explicit record-to-polynomial radix construction, with ordered
recovery, exact capacity \(t\ge B^\alpha\), zero padding and a native
interface retaining XPIR's linear reply arithmetic.
\item Scoped formal analysis: conditional integer no-wrap correctness,
a parametric full-task bound for a specified rounded Gaussian law, and
single-challenge index privacy from uniform-secret decisional PLWE with
\(N+1\) samples and explicit unit-probability loss.
\item Separate controlled evaluations of two-record task latency,
multiplicity, implementation/resource sensitivity and capacity, together
with reproducible native-buffer accounting and functional checks.
\end{itemize}

The primary two-record comparison covers 48 workload/profile points in
both cold and online views. Its median paired repeated-to-packed latency
ratios range from 1.059 to 1.190, with all 96 pointwise intervals above
one. Further experiments extend multiplicity and block coverage under
their own contracts; they are not pooled with the primary evidence.
These findings concern a fixed profile and host. They establish neither
a global XPIR optimum nor superiority over other PIR systems, and the
native/formal sampler and reliability bridge remains unresolved.
